\documentclass{article}
\usepackage[left=2cm, right=2cm, top=2cm, bottom=2cm]{geometry}
\usepackage{graphicx}
\usepackage{amsmath}
\usepackage{amssymb}
\usepackage{amsthm}
\usepackage{fancyhdr}
\usepackage{verbatim}
\usepackage{listings}
\usepackage{xcolor}
\usepackage{pdfpages}
\usepackage{cancel}
\usepackage{physics}
\usepackage{esint}
\usepackage{braket}

\lstset{
    language=C++,
    basicstyle=\ttfamily\footnotesize,
    keywordstyle=\color{blue},
    commentstyle=\color{green},
    stringstyle=\color{red},
    numbers=left,
    numberstyle=\tiny\color{gray},
    frame=single,
    breaklines=true,
    captionpos=b,
}


\title{MTH 446: Lecture \# 11}
\author{Cliff Sun}

\newtheorem{theorem}{Theorem}[section]
\newtheorem{lemma}[theorem]{Lemma}
\newtheorem{definition}[theorem]{Definition}
\newtheorem{conjecture}[theorem]{Conjecture}
\newtheorem{proposition}[theorem]{Proposition}
\newtheorem{corollary}[theorem]{Corollary}
\newtheorem{one minute paper}[theorem]{One Minute Paper}

\pagestyle{fancy}
\lhead{\textbf{\thepage}\ \ \nouppercase{\rightmark}}
\chead{MTH 446: Lecture \# 11}
\rhead{Cliff Sun}

\begin{document}

\maketitle

\section*{Lecture Span}
\begin{enumerate}
    \item CR equations
\end{enumerate}

\section*{Cauchy Riemann Equations}

The derivative w.r.t the real number line is the same as that on the complex number line. Note that $z = z_0 + h$, then 

\begin{align}
    x + iy &= (x_0 + h) + iy_0
\end{align}

Therefore, 

\begin{align}
    \frac{f(z) - f(z_0)}{z - z_0} &= \frac{(u(x + h, y) - u(x,y)) + i(v(x + h, y) - v(x,y))}{h} \\
    &= \partial_x u + i \partial_x v
\end{align}

Similarly, taking the partial derivative on the complex plane yields:

\begin{align}
    \frac{f(z) - f(z_0)}{z - z_0} &= \frac{(u(x, y + h) - u(x,y)) + i(v(x, y + h) - v(x,y))}{ih} \\
    &= -i\partial_y u + \partial_y v
\end{align}

Therefore, 

\begin{equation}
    \partial_x u = \partial_y v, \quad \partial_x v = -\partial_y u
\end{equation}

\subsection*{Example}

Let $f(z) = x + iy$. Here, the CR equations are not satisfied. Therefore, $f'(z)$ DNE for all $z \in \mathbb{C}$. 

\subsection*{Example}

Consider $f(z) = x^2 + y^2$. $f'(z)$ DNE if $z \neq 0$ by the necessary conditions of complex differentiability. next, 

\begin{equation}
    \lim_{z \rightarrow 0}\frac{|z|^2}{z} = \lim \overline{z} = 0
\end{equation}

Therefore, $f$ is complex differentiable only at the origin. 

\subsection*{Example}

$f(z) = \overline{z}\Im(z)$. Here, 

\begin{equation}
    f(z) = xy - iy^2
\end{equation}

Here, we obtain that $f(z)$ is not complex differentiable for $z \neq 0$ by the necessary condition of complex differentiability. 

At $f'(z)$, 

\begin{equation}
    f'(z) = \lim \frac{\overline{z}\Im(z)}{z} = \lim \frac{|\overline{z}| |\Im(z)|}{|z|} = \lim |y| = 0
\end{equation}

Therefore, $f$ is complex differentiable only at the origin. 

\subsection*{Example}

For $f(z) = z^2 = (x^2 + y^2) + 2xy i$. The C.R. equations are satisfied at every point, but we cannot conclude that $d/dz (z^2)$ exists at every point $z \in \mathbb{C}$. We need the next theorem... 

\begin{theorem}
    \underline{Sufficient conditions for complex differentiability}. Let $f: G \rightarrow \mathbb{C}$ be a function. If $u_x,u_y,v_x,v_y$ are continuous at $z = z_0$ and the CR equations 
    are satisfied at $z = z_0$. Then $f'(z_0)$ exists.
\end{theorem}

\subsection*{Example}

Use the previous example $f(z) = z^2$. Then, the CR equations are satisfied everywhere, and thus $f$ is CD everywhere on $\mathbb{C}$. 

\textbf{Must verify that $v_i$ and $u_i$ are continuous!!!}

\subsection*{Example}

$f(z) = 2x + ixy$. 

Formal ending: ''By the necessary condition of complex differentiability, $f'(z)$ does not exist if $z \neq 2$". 

\end{document}